% ============================================================
%  RÉSUMÉ (Français)
% ============================================================
\chapter*{Résumé}
\addcontentsline{toc}{chapter}{Résumé}
\thispagestyle{plain}

Ce document est le résultat d'un projet de fin d'études réalisé à l'École Hassania des Travaux Publics (EHTP), dans le cadre de l'obtention du diplôme d'Ingénieur d'État en Génie de l'Eau et de l'Environnement.

L'objectif principal de ce projet est de simuler numériquement les écoulements hydrauliques au sein des ouvrages du \textbf{barrage Ahl Souss (Ait M'zal)}, en particulier l'évacuateur de crues à seuil libre de type Creager et la vidange de fond. Cette approche s'appuie sur la \textbf{Mécanique des Fluides Numérique (CFD)} via le logiciel libre \textbf{OpenFOAM}, en utilisant la méthode des \textbf{volumes finis} pour la discrétisation des équations de Navier-Stokes et le modèle de turbulence \textbf{$k$-$\omega$ SST}.

La démarche adoptée comprend une étude bibliographique approfondie des fondements de la CFD, des modèles de turbulence et des méthodes de discrétisation numérique, suivie d'une présentation détaillée du barrage. Les géométries des ouvrages ont été construites sous \textbf{Autodesk Revit} puis maillées pour la simulation. Les résultats numériques (champs de vitesse, pression, lignes de courant) sont comparés aux données théoriques et dimensionnelles disponibles.

\vspace{0.5cm}
\noindent\hrulefill

\noindent\textbf{Mots clés :}\quad CFD, OpenFOAM, hydraulique, barrage, évacuateur de crues, vidange de fond, volumes finis, turbulence, $k$-$\omega$ SST, Navier-Stokes.

\noindent\hrulefill
\vspace{0.3cm}

\newpage

% ============================================================
%  ABSTRACT (English)
% ============================================================
\chapter*{Abstract}
\addcontentsline{toc}{chapter}{Abstract}
\thispagestyle{plain}

This document is the result of an end-of-studies project carried out at the École Hassania des Travaux Publics (EHTP), as part of the requirements for obtaining the State Engineer's degree in Water and Environmental Engineering.

The main objective of this project is to numerically simulate hydraulic flows within the hydraulic structures of the \textbf{Ahl Souss (Ait M'zal) dam}, specifically the Creager-type free-overflow spillway and the bottom outlet. This approach relies on \textbf{Computational Fluid Dynamics (CFD)} using the open-source software \textbf{OpenFOAM}, employing the \textbf{finite volume method} to discretize the Navier-Stokes equations and the \textbf{$k$-$\omega$ SST} turbulence model.

The methodology includes a comprehensive literature review of CFD fundamentals, turbulence models, and numerical discretization methods, followed by a detailed presentation of the dam. The geometries of the structures were built using \textbf{Autodesk Revit} and then meshed for simulation. Numerical results (velocity fields, pressure, streamlines) are compared with theoretical and dimensional data.

\vspace{0.5cm}
\noindent\hrulefill

\noindent\textbf{Keywords:}\quad CFD, OpenFOAM, hydraulics, dam, spillway, bottom outlet, finite volumes, turbulence, $k$-$\omega$ SST, Navier-Stokes.

\noindent\hrulefill
\vspace{0.3cm}

\newpage
